% ==============================================================================
% CỤM 2 - ĐỀ THAM KHẢO 2 (NĂM HỌC 2026 - 2027)
% ==============================================================================

\begin{bt}%[9D6H18-2]
% nguon: Đề tham khảo Tuyển sinh 10 TP.HCM 2026-2027 – Cụm 2 (Đề 2), Bài 1, trang 11
Cho hàm số $(P): y = \dfrac{1}{2}x^2$.
\begin{enumerate}[a)]
	\item Vẽ đồ thị $(P)$.
	\item Tìm điểm $M$ thuộc $(P)$ có tung độ gấp đôi hoành độ và khác gốc tọa độ.
\end{enumerate}
\loigiai{
	\begin{enumerate}[a)]
		\item Bảng giá trị tương ứng giữa $x$ và $y$:
		\begin{center}
		\begin{tabular}{|c|c|c|c|c|c|}
		\hline
		$x$ & $-4$ & $-2$ & $0$ & $2$ & $4$ \\
		\hline
		$y = \dfrac{1}{2}x^2$ & $8$ & $2$ & $0$ & $2$ & $8$ \\
		\hline
		\end{tabular}
		\end{center}
		Đồ thị $(P)$ là đường cong parabol đi qua gốc tọa độ $O(0; 0)$ và các điểm $(-4; 8)$, $(-2; 2)$, $(2; 2)$, $(4; 8)$, nhận trục tung $Oy$ làm trục đối xứng và có bề lõm quay lên trên.
		\begin{center}
		\begin{tikzpicture}[>=stealth, scale=0.7]
			\draw[->] (-5,0) -- (5,0) node[below] {$x$};
			\draw[->] (0,-1) -- (0,9) node[left] {$y$};
			\draw (0,0) node[below left] {$O$};
			\foreach \x in {-4,-2,2,4}
				\draw (\x,0.1) -- (\x,-0.1) node[below] {\small $\x$};
			\foreach \y in {2,8}
				\draw (0.1,\y) -- (-0.1,\y) node[left] {\small $\y$};
			\draw[dashed] (-4,0) -- (-4,8) -- (4,8) -- (4,0);
			\draw[dashed] (-2,0) -- (-2,2) -- (2,2) -- (2,0);
			\draw[smooth,samples=100,domain=-4.2:4.2,thick,blue] plot(\x,{0.5*(\x)^2});
			\node[right] at (3.5,6.5) {$(P): y=\dfrac{1}{2}x^2$};
			\fill (-4,8) circle (1.5pt) (-2,2) circle (1.5pt) (0,0) circle (1.5pt) (2,2) circle (1.5pt) (4,8) circle (1.5pt);
		\end{tikzpicture}
		\end{center}
		\item Gọi $M(x; y) \in (P)$. Vì tung độ gấp đôi hoành độ nên:
		$$y = 2x.$$
		Thay vào phương trình parabol $(P)$, ta có:
		$$\frac{1}{2}x^2 = 2x \iff \frac{1}{2}x^2 - 2x = 0 \iff x\left(\frac{1}{2}x - 2\right) = 0 \iff \left[\begin{aligned} x &= 0 \\ x &= 4. \end{aligned}\right.$$
		Vì điểm $M$ khác gốc tọa độ $O(0; 0)$ nên ta nhận $x = 4$.\\
		Với $x = 4 \implies y = 2 \cdot 4 = 8$.\\
		Vậy tọa độ điểm cần tìm là $M(4; 8)$.
	\end{enumerate}
}
\end{bt}

\begin{bt}%[9D6H20-2]
% nguon: Đề tham khảo Tuyển sinh 10 TP.HCM 2026-2027 – Cụm 2 (Đề 2), Bài 2, trang 11
Cho phương trình: $2x^2 - 8x - 3 = 0$.
\begin{enumerate}[a)]
	\item Chứng minh phương trình trên luôn có hai nghiệm phân biệt.
	\item Không giải phương trình, hãy tính giá trị biểu thức:
	$$A = x_1(x_1 - 2) + x_2(x_2 - 2)$$
\end{enumerate}
\loigiai{
	\begin{enumerate}[a)]
		\item Xét phương trình $2x^2 - 8x - 3 = 0$ có $a = 2, b' = -4, c = -3$.\\
		Vì $a \cdot c = 2 \cdot (-3) = -6 < 0$ nên phương trình luôn có hai nghiệm phân biệt trái dấu $x_1, x_2$.\\
		\textit{(Cách khác: $\Delta' = (-4)^2 - 2 \cdot (-3) = 16 + 6 = 22 > 0$ nên phương trình luôn có hai nghiệm phân biệt).}
		\item Theo định lý Vi-ét, ta có:
		$$\begin{cases} x_1 + x_2 = -\dfrac{b}{a} = -\dfrac{-8}{2} = 4 \\ x_1 x_2 = \dfrac{c}{a} = -\dfrac{3}{2} \end{cases}$$
		Từ đó tính được:
		$$x_1^2 + x_2^2 = (x_1 + x_2)^2 - 2x_1 x_2 = 4^2 - 2 \cdot \left(-\frac{3}{2}\right) = 16 + 3 = 19.$$
		Biến đổi biểu thức $A$:
		$$A = x_1(x_1 - 2) + x_2(x_2 - 2) = x_1^2 - 2x_1 + x_2^2 - 2x_2 = (x_1^2 + x_2^2) - 2(x_1 + x_2).$$
		Thay các giá trị tương ứng vào biểu thức $A$, ta được:
		$$A = 19 - 2 \cdot 4 = 19 - 8 = 11.$$
		Vậy $A = 11$.
	\end{enumerate}
}
\end{bt}

\begin{bt}%[9D8H26-3]
% nguon: Đề tham khảo Tuyển sinh 10 TP.HCM 2026-2027 – Cụm 2 (Đề 2), Bài 3, trang 11
Số lượng máy tính để bàn và máy tính xách tay được bán ra trong $6$ tháng đầu năm của một công ty được ghi nhận như sau (đơn vị: chiếc máy):
\begin{itemize}
	\item Tháng Một: $80$ máy để bàn, $30$ máy xách tay.
	\item Tháng Hai: $41$ máy để bàn, $21$ máy xách tay.
	\item Tháng Ba: $23$ máy để bàn, $38$ máy xách tay.
	\item Tháng Bốn: $20$ máy để bàn, $50$ máy xách tay.
	\item Tháng Năm: $14$ máy để bàn, $43$ máy xách tay.
	\item Tháng Sáu: $7$ máy để bàn, $59$ máy xách tay.
\end{itemize}
\begin{enumerate}[a)]
	\item Trong $6$ tháng đầu năm, tháng nào có sự chênh lệch giữa số lượng máy tính xách tay và máy tính để bàn bán ra là ít nhất?
	\item Chọn ngẫu nhiên $1$ tháng trong $6$ tháng đầu năm, tính xác suất của các biến cố sau:
	\begin{itemize}
		\item Biến cố $A$: "Tháng được chọn có số lượng máy tính để bàn mà công ty bán được không quá $30$ máy".
		\item Biến cố $B$: "Tháng được chọn có sự chênh lệch số lượng giữa hai loại máy bán ra không quá $20$ máy".
	\end{itemize}
\end{enumerate}
\loigiai{
	\begin{enumerate}[a)]
		\item Độ chênh lệch giữa số lượng máy tính xách tay và máy tính để bàn trong từng tháng là:
		\begin{itemize}
			\item Tháng Một: $|80 - 30| = 50$ (máy).
			\item Tháng Hai: $|41 - 21| = 20$ (máy).
			\item Tháng Ba: $|38 - 23| = 15$ (máy).
			\item Tháng Bốn: $|50 - 20| = 30$ (máy).
			\item Tháng Năm: $|43 - 14| = 29$ (máy).
			\item Tháng Sáu: $|59 - 7| = 52$ (máy).
		\end{itemize}
		So sánh các kết quả trên, ta thấy $15 < 20 < 29 < 30 < 50 < 52$.\\
		Vậy tháng Ba có sự chênh lệch giữa số lượng hai loại máy bán ra là ít nhất ($15$ máy).
		\item Không gian mẫu là tập hợp $6$ tháng đầu năm:
		$$\Omega = \{\text{Tháng Một}; \text{Tháng Hai}; \text{Tháng Ba}; \text{Tháng Bốn}; \text{Tháng Năm}; \text{Tháng Sáu}\}.$$
		Số phần tử của không gian mẫu là $n(\Omega) = 6$.
		\begin{itemize}
			\item Biến cố $A$: "Tháng được chọn có số lượng máy tính để bàn không quá $30$ máy".\\
			Các tháng có số máy để bàn $\le 30$ là: Tháng Ba ($23$ máy), Tháng Bốn ($20$ máy), Tháng Năm ($14$ máy), Tháng Sáu ($7$ máy).\\
			Số kết quả thuận lợi cho $A$ là $n(A) = 4$.\\
			Xác suất của biến cố $A$ là:
			$$P(A) = \frac{n(A)}{n(\Omega)} = \frac{4}{6} = \frac{2}{3}.$$
			\item Biến cố $B$: "Tháng được chọn có sự chênh lệch số lượng giữa hai loại máy không quá $20$ máy".\\
			Các tháng có mức chênh lệch $\le 20$ máy là: Tháng Hai ($20$ máy) và Tháng Ba ($15$ máy).\\
			Số kết quả thuận lợi cho $B$ là $n(B) = 2$.\\
			Xác suất của biến cố $B$ là:
			$$P(B) = \frac{n(B)}{n(\Omega)} = \frac{2}{6} = \frac{1}{3}.$$
		\end{itemize}
	\end{enumerate}
}
\end{bt}

\begin{bt}%[9D6V21-3]
% nguon: Đề tham khảo Tuyển sinh 10 TP.HCM 2026-2027 – Cụm 2 (Đề 2), Bài 4, trang 11
Tại một giao lộ, người ta thiết kế một vòng xuyến gồm một hình tròn có bán kính $a\text{ (m)}$ trồng hoa hồng và một hình vành khuyên có độ rộng $a - 1\text{ (m)}$ bao quanh hình tròn đó để trồng hoa cúc ($a > 1$).
\begin{enumerate}[a)]
	\item Viết biểu thức $M$ theo biến $a$ biểu diễn diện tích phần đất trồng hoa cúc.
	\item Biết rằng diện tích phần đất trồng hoa hồng bằng với diện tích phần đất trồng hoa cúc. Hãy tính giá trị của $a$ (làm tròn kết quả đến hàng phần mười).
\end{enumerate}
\loigiai{
	\begin{enumerate}[a)]
		\item Bán kính của hình tròn trồng hoa hồng là $R_1 = a\text{ (m)}$.\\
		Bán kính vòng ngoài của hình vành khuyên là:
		$$R_2 = a + (a - 1) = 2a - 1\text{ (m)}.$$
		Diện tích của hình tròn lớn (vòng ngoài) là:
		$$S_{\text{ngoài}} = \pi R_2^2 = \pi(2a - 1)^2\text{ (m}^2\text{)}.$$
		Diện tích hình tròn trồng hoa hồng là:
		$$S_{\text{hồng}} = \pi R_1^2 = \pi a^2\text{ (m}^2\text{)}.$$
		Diện tích phần đất hình vành khuyên trồng hoa cúc là:
		$$M = S_{\text{ngoài}} - S_{\text{hồng}} = \pi(2a - 1)^2 - \pi a^2 = \pi(4a^2 - 4a + 1 - a^2) = \pi(3a^2 - 4a + 1)\text{ (m}^2\text{)}.$$
		\item Theo đề bài, diện tích trồng hoa hồng bằng diện tích trồng hoa cúc nên:
		$$\pi a^2 = \pi(3a^2 - 4a + 1) \iff a^2 = 3a^2 - 4a + 1 \iff 2a^2 - 4a + 1 = 0.$$
		Phương trình có $\Delta' = (-2)^2 - 2 \cdot 1 = 4 - 2 = 2 > 0$.\\
		Suy ra phương trình có hai nghiệm:
		$$a_1 = \frac{2 + \sqrt{2}}{2} \approx 1{,}7, \quad a_2 = \frac{2 - \sqrt{2}}{2} \approx 0{,}3.$$
		Vì điều kiện $a > 1$ nên ta nhận $a = \dfrac{2 + \sqrt{2}}{2} \approx 1{,}7\text{ (m)}$.\\
		Vậy giá trị của $a$ xấp xỉ bằng $1{,}7\text{ m}$.
	\end{enumerate}
}
\end{bt}

\begin{bt}%[9H0V31-3]
% nguon: Đề tham khảo Tuyển sinh 10 TP.HCM 2026-2027 – Cụm 2 (Đề 2), Bài 5, trang 12
Gạch chống nóng $6$ lỗ (gạch Tuynel) có dạng hình hộp chữ nhật với kích thước dài $195\text{ mm}$, rộng $135\text{ mm}$ và cao $90\text{ mm}$. Mỗi viên gạch có $6$ lỗ rỗng hình trụ bằng nhau chạy suốt chiều dài viên gạch, mỗi lỗ có đường kính đáy là $28\text{ mm}$.
\begin{enumerate}[a)]
	\item Lấy $\pi = 3{,}14$. Hãy tính thể tích nguyên vật liệu đất sét để làm nên một viên gạch trên (bỏ qua các rãnh gân của viên gạch).
	\item Một khối đất nung dạng hình hộp chữ nhật có kích thước $2{,}1\text{ m} \times 1{,}5\text{ m} \times 1{,}5\text{ m}$. Người ta dùng khối đất đó để làm gạch. Hỏi cần ít nhất bao nhiêu khối đất như trên (lấy số nguyên) để làm ra $10\,000$ viên gạch, biết hao hụt đất nung trong quá trình làm gạch là $10\%$? Cho biết công thức thể tích hình trụ là $V = \pi r^2 h$, hình hộp chữ nhật là $V = abc$.
\end{enumerate}
\loigiai{
	\begin{enumerate}[a)]
		\item Thể tích toàn phần của khối hình hộp chữ nhật bên ngoài viên gạch là:
		$$V_{\text{hộp}} = 195 \cdot 135 \cdot 90 = 2\,369\,250\text{ (mm}^3\text{)}.$$
		Bán kính đáy của mỗi lỗ rỗng hình trụ là:
		$$r = \frac{28}{2} = 14\text{ (mm)}.$$
		Chiều cao của mỗi lỗ rỗng bằng chiều dài viên gạch là $h = 195\text{ mm}$.\\
		Tổng thể tích của $6$ lỗ rỗng hình trụ là:
		$$V_{\text{lỗ}} = 6 \cdot (\pi r^2 h) = 6 \cdot 3{,}14 \cdot 14^2 \cdot 195 = 719\,740{,}8\text{ (mm}^3\text{)}.$$
		Thể tích nguyên vật liệu đất sét thực tế để làm một viên gạch là:
		$$V = V_{\text{hộp}} - V_{\text{lỗ}} = 2\,369\,250 - 719\,740{,}8 = 1\,649\,509{,}2\text{ (mm}^3\text{)}.$$
		\item Tổng thể tích đất sét cần cho $10\,000$ viên gạch thành phẩm là:
		$$V_{10\,000} = 10\,000 \cdot 1\,649\,509{,}2 = 16\,495\,092\,000\text{ (mm}^3\text{)} = 16{,}495092\text{ (m}^3\text{)}.$$
		Vì hao hụt đất nung trong quá trình sản xuất là $10\%$ nên lượng đất sét thô thực tế cần chuẩn bị là:
		$$V_{\text{cần}} = \frac{V_{10\,000}}{100\% - 10\%} = \frac{16{,}495092}{0{,}9} = 18{,}32788\text{ (m}^3\text{)}.$$
		Thể tích của một khối đất nung hình hộp chữ nhật là:
		$$V_{\text{khối}} = 2{,}1 \cdot 1{,}5 \cdot 1{,}5 = 4{,}725\text{ (m}^3\text{)}.$$
		Số khối đất nung cần chuẩn bị là:
		$$\frac{V_{\text{cần}}}{V_{\text{khối}}} = \frac{18{,}32788}{4{,}725} \approx 3{,}88\text{ (khối)}.$$
		Vì số khối đất phải là số nguyên nên cần ít nhất $4$ khối đất nung như trên.
	\end{enumerate}
}
\end{bt}

\begin{bt}%[9D1H3-4]
% nguon: Đề tham khảo Tuyển sinh 10 TP.HCM 2026-2027 – Cụm 2 (Đề 2), Bài 6, trang 12
Một hôm hai bố con đang đi dạo trong một dãy phố có $6$ căn nhà liên tiếp nằm cùng một phía bên phải của con đường, được đánh số địa chỉ là các số chẵn liên tiếp. Người bố đố con: "Tìm thương và số dư khi lấy số địa chỉ nhà cuối dãy chia cho số địa chỉ nhà đầu dãy", thì em bé trả lời thương là $2$ và số dư là $2$. Biết rằng câu trả lời của em bé là hoàn toàn chính xác. Hãy xác định số địa chỉ của từng căn nhà trong dãy phố đó.
\loigiai{
	Gọi số địa chỉ của căn nhà đầu dãy là $x$ và nhà cuối dãy là $y$ ($x, y \in \mathbb{N}^*$, $x, y$ là các số chẵn).\\
	Vì dãy phố có $6$ căn nhà liên tiếp đánh số chẵn liên tiếp nên hiệu giữa số nhà cuối và số nhà đầu là:
	$$y - x = (6 - 1) \cdot 2 = 10 \quad (1)$$
	Theo đề bài, khi lấy số nhà cuối chia cho số nhà đầu được thương là $2$ và số dư là $2$, ta có:
	$$y = 2x + 2 \quad (2)$$
	Thay $(2)$ vào $(1)$, ta được:
	$$(2x + 2) - x = 10 \iff x + 2 = 10 \iff x = 8\text{ (thỏa mãn).}$$
	Từ đó tính được số nhà cuối dãy:
	$$y = 8 + 10 = 18.$$
	Vậy số địa chỉ nhà đầu tiên là $8$ và các số nhà tiếp theo là các số chẵn liên tiếp.\\
	Địa chỉ của $6$ căn nhà lần lượt là: $8$, $10$, $12$, $14$, $16$, $18$.
}
\end{bt}

\begin{bt}%[9H9C29-3]
% nguon: Đề tham khảo Tuyển sinh 10 TP.HCM 2026-2027 – Cụm 2 (Đề 2), Bài 7, trang 12
Cho đường tròn $(O)$ có đường kính $BC$. Trên nửa đường tròn $(O)$, lấy hai điểm $A$ và $D$ sao cho số đo cung $BA$ nhỏ hơn số đo cung $BD$. Tia $BA$ và tia $CD$ cắt nhau tại $S$, đoạn thẳng $AC$ cắt $BD$ tại $H$. Gọi $I$ là trung điểm của $SH$.
\begin{enumerate}[a)]
	\item Tính số đo $\widehat{BAC}$ và chứng minh tứ giác $SAHD$ nội tiếp.
	\item Tia $SH$ cắt $BC$ tại $M$. Chứng minh $\widehat{IAH} = \widehat{MDC}$.
	\item Trong trường hợp $\widehat{BSC} = 60^\circ$ và $BC = 6\text{ cm}$, hãy tính độ dài đoạn thẳng $AD$ và bán kính đường tròn ngoại tiếp tam giác $SAD$.
\end{enumerate}
\loigiai{
	\begin{center}
	\begin{tikzpicture}[scale=1]
		\coordinate (O) at (0,0);
		\def\R{3}
		\coordinate (B) at (-\R,0);
		\coordinate (C) at (\R,0);
		\draw[thick] (B) -- (C);
		\draw[thick] (C) arc (0:180:\R);
		
		\coordinate (A) at (140:\R);
		\coordinate (D) at (50:\R);
		
		\coordinate (S) at (intersection of B--A and C--D);
		\coordinate (H) at (intersection of A--C and B--D);
		\coordinate (I) at ($(S)!0.5!(H)$);
		\coordinate (M) at ($(B)!(S)!(C)$);
		
		\draw[thick] (B) -- (S) -- (C);
		\draw[thick] (A) -- (C);
		\draw[thick] (B) -- (D);
		\draw[thick, blue] (S) -- (M);
		\draw[thick, red] (A) -- (D);
		\draw[dashed] (I) -- (A);
		\draw[dashed] (M) -- (D);
		
		\fill (O) circle (1.5pt) node[below] {$O$};
		\fill (B) circle (1.5pt) node[below left] {$B$};
		\fill (C) circle (1.5pt) node[below right] {$C$};
		\fill (A) circle (1.5pt) node[above left] {$A$};
		\fill (D) circle (1.5pt) node[above right] {$D$};
		\fill (S) circle (1.5pt) node[above] {$S$};
		\fill (H) circle (1.5pt) node[below] {$H$};
		\fill (I) circle (1.5pt) node[left] {$I$};
		\fill (M) circle (1.5pt) node[below] {$M$};
	\end{tikzpicture}
	\end{center}
	\begin{enumerate}[a)]
		\item Ta có $\widehat{BAC} = 90^\circ$ (góc nội tiếp chắn nửa đường tròn đường kính $BC$).\\
		Tương tự, $\widehat{BDC} = 90^\circ$ (góc nội tiếp chắn nửa đường tròn đường kính $BC$).\\
		Từ đó suy ra:
		$$\widehat{SAH} = 180^\circ - \widehat{BAC} = 90^\circ \quad \text{và} \quad \widehat{SDH} = 180^\circ - \widehat{BDC} = 90^\circ.$$
		Xét tứ giác $SAHD$ có $\widehat{SAH} + \widehat{SDH} = 90^\circ + 90^\circ = 180^\circ$, do đó tứ giác $SAHD$ nội tiếp đường tròn đường kính $SH$.
		\item Vì tứ giác $SAHD$ nội tiếp đường tròn đường kính $SH$ có tâm là $I$ (trung điểm của $SH$) nên:
		$$IA = IH = R_I \implies \Delta IAH \text{ cân tại } I \implies \widehat{IAH} = \widehat{IHA}.$$
		Mà $\widehat{IHA} = \widehat{MHC}$ (hai góc đối đỉnh).\\
		Trong tam giác $SBC$, ta có $CA \perp SB$ (do $\widehat{BAC} = 90^\circ$) và $BD \perp SC$ (do $\widehat{BDC} = 90^\circ$), nên $CA$ và $BD$ là hai đường cao cắt nhau tại $H$. Suy ra $H$ là trực tâm của tam giác $SBC$.\\
		Do đó $SH \perp BC$ tại $M$, hay tam giác $HMC$ vuông tại $M$.\\
		Xét tứ giác $DMCH$ có:
		$$\widehat{HDC} = 90^\circ \quad \text{và} \quad \widehat{HMC} = 90^\circ.$$
		Hai đỉnh $D$ và $M$ cùng nhìn đoạn thẳng $HC$ dưới một góc vuông, nên tứ giác $DMCH$ nội tiếp đường tròn đường kính $HC$.\\
		Suy ra $\widehat{MDC} = \widehat{MHC}$ (hai góc nội tiếp cùng chắn cung $MC$).\\
		Từ đó ta có $\widehat{IAH} = \widehat{MDC}$.
		\item Xét $\Delta SAD$ và $\Delta SCB$ có:
		$$\widehat{BSC}\text{ là góc chung và } \widehat{SAD} = \widehat{SCB}\text{ (cùng bù với } \widehat{BAD}\text{)}.$$
		Suy ra $\Delta SAD \backsim \Delta SCB$ (g-g), do đó:
		$$\frac{AD}{BC} = \frac{SA}{SC} = \cos \widehat{BSC} = \cos 60^\circ = \frac{1}{2}.$$
		Suy ra:
		$$AD = \frac{1}{2}BC = \frac{1}{2} \cdot 6 = 3\text{ (cm)}.$$
		Đường tròn ngoại tiếp tam giác $SAD$ chính là đường tròn $(I)$ ngoại tiếp tứ giác $SAHD$.\\
		Trong đường tròn $(I)$, ta có góc ở tâm:
		$$\widehat{AID} = 2\widehat{ASD} = 2\widehat{BSC} = 2 \cdot 60^\circ = 120^\circ.$$
		Kẻ $IQ \perp AD$ tại $Q$. Do $\Delta IAD$ cân tại $I$ ($IA = ID = R_I$) nên $IQ$ đồng thời là đường phân giác và đường trung tuyến ứng với cạnh đáy $AD$.\\
		Suy ra $QD = \dfrac{1}{2}AD = 1{,}5\text{ cm}$ và $\widehat{QID} = \dfrac{1}{2}\widehat{AID} = 60^\circ$.\\
		Xét tam giác vuông $IQD$ vuông tại $Q$:
		$$\sin \widehat{QID} = \frac{QD}{ID} \implies ID = \frac{QD}{\sin 60^\circ} = \frac{1{,}5}{\frac{\sqrt{3}}{2}} = \frac{3}{\sqrt{3}} = \sqrt{3}\text{ (cm)}.$$
		Vậy $AD = 3\text{ cm}$ và bán kính đường tròn ngoại tiếp tam giác $SAD$ là $\sqrt{3}\text{ cm}$.
	\end{enumerate}
}
\end{bt}

% ==============================================================================
% CỤM 2 - ĐỀ THAM KHẢO 3 (NĂM HỌC 2026 - 2027)
% ==============================================================================

\begin{bt}%[9D6H18-2]
% nguon: Đề tham khảo Tuyển sinh 10 TP.HCM 2026-2027 – Cụm 2 (Đề 3), Bài 1, trang 13
Cho hàm số $(P): y = -\dfrac{1}{2}x^2$.
\begin{enumerate}[a)]
	\item Vẽ đồ thị $(P)$ của hàm số trên.
	\item Tìm những điểm $A$ thuộc $(P)$ sao cho tung độ bằng $-\dfrac{1}{2}$.
\end{enumerate}
\loigiai{
	\begin{enumerate}[a)]
		\item Bảng giá trị tương ứng giữa $x$ và $y$:
		\begin{center}
		\begin{tabular}{|c|c|c|c|c|c|}
		\hline
		$x$ & $-4$ & $-2$ & $0$ & $2$ & $4$ \\
		\hline
		$y = -\dfrac{1}{2}x^2$ & $-8$ & $-2$ & $0$ & $-2$ & $-8$ \\
		\hline
		\end{tabular}
		\end{center}
		Đồ thị $(P)$ là đường cong parabol đi qua gốc tọa độ $O(0; 0)$ và các điểm $(-4; -8)$, $(-2; -2)$, $(2; -2)$, $(4; -8)$, nhận trục tung $Oy$ làm trục đối xứng và có bề lõm quay xuống dưới.
		\begin{center}
		\begin{tikzpicture}[>=stealth, scale=0.7]
			\draw[->] (-5,0) -- (5,0) node[below] {$x$};
			\draw[->] (0,-9) -- (0,1.5) node[left] {$y$};
			\draw (0,0) node[above right] {$O$};
			\foreach \x in {-4,-2,2,4}
				\draw (\x,0.1) -- (\x,-0.1) node[above] {\small $\x$};
			\foreach \y in {-2,-8}
				\draw (0.1,\y) -- (-0.1,\y) node[left] {\small $\y$};
			\draw[dashed] (-4,0) -- (-4,-8) -- (4,-8) -- (4,0);
			\draw[dashed] (-2,0) -- (-2,-2) -- (2,-2) -- (2,0);
			\draw[smooth,samples=100,domain=-4.2:4.2,thick,blue] plot(\x,{-0.5*(\x)^2});
			\node[right] at (3.5,-4.5) {$(P): y=-\dfrac{1}{2}x^2$};
			\fill (-4,-8) circle (1.5pt) (-2,-2) circle (1.5pt) (0,0) circle (1.5pt) (2,-2) circle (1.5pt) (4,-8) circle (1.5pt);
		\end{tikzpicture}
		\end{center}
		\item Điểm $A$ thuộc $(P)$ có tung độ bằng $-\dfrac{1}{2}$ nên $y = -\dfrac{1}{2}$.\\
		Thay $y = -\dfrac{1}{2}$ vào phương trình parabol $(P)$, ta có:
		$$-\frac{1}{2}x^2 = -\frac{1}{2} \iff x^2 = 1 \iff x = 1 \text{ hoặc } x = -1.$$
		Với $x = 1 \implies y = -\dfrac{1}{2}$, ta được điểm $A_1\left(1; -\dfrac{1}{2}\right)$.\\
		Với $x = -1 \implies y = -\dfrac{1}{2}$, ta được điểm $A_2\left(-1; -\dfrac{1}{2}\right)$.\\
		Vậy có hai điểm thỏa mãn yêu cầu là $\left(1; -\dfrac{1}{2}\right)$ và $\left(-1; -\dfrac{1}{2}\right)$.
	\end{enumerate}
}
\end{bt}

\begin{bt}%[9D6H20-2]
% nguon: Đề tham khảo Tuyển sinh 10 TP.HCM 2026-2027 – Cụm 2 (Đề 3), Bài 2, trang 13
Cho phương trình: $3x^2 - 2x - 5 = 0$.
\begin{enumerate}[a)]
	\item Chứng tỏ phương trình luôn có hai nghiệm phân biệt $x_1, x_2$.
	\item Không giải phương trình, hãy tính giá trị của biểu thức:
	$$M = x_1(x_1 - 1) + x_2(x_2 - 1)$$
\end{enumerate}
\loigiai{
	\begin{enumerate}[a)]
		\item Xét phương trình $3x^2 - 2x - 5 = 0$ có các hệ số $a = 3, b = -2, c = -5$.\\
		Vì $a \cdot c = 3 \cdot (-5) = -15 < 0$ nên phương trình luôn có hai nghiệm phân biệt trái dấu $x_1, x_2$.\\
		\textit{(Cách khác: $\Delta' = (-1)^2 - 3 \cdot (-5) = 1 + 15 = 16 > 0$ nên phương trình luôn có hai nghiệm phân biệt).}
		\item Theo định lý Vi-ét, ta có:
		$$\begin{cases} x_1 + x_2 = -\dfrac{b}{a} = \dfrac{2}{3} \\ x_1 x_2 = \dfrac{c}{a} = -\dfrac{5}{3} \end{cases}$$
		Từ đó tính được:
		$$x_1^2 + x_2^2 = (x_1 + x_2)^2 - 2x_1 x_2 = \left(\frac{2}{3}\right)^2 - 2 \cdot \left(-\frac{5}{3}\right) = \frac{4}{9} + \frac{10}{3} = \frac{34}{9}.$$
		Biến đổi biểu thức $M$:
		$$M = x_1(x_1 - 1) + x_2(x_2 - 1) = x_1^2 - x_1 + x_2^2 - x_2 = (x_1^2 + x_2^2) - (x_1 + x_2).$$
		Thay các giá trị tương ứng vào biểu thức $M$, ta được:
		$$M = \frac{34}{9} - \frac{2}{3} = \frac{34 - 6}{9} = \frac{28}{9}.$$
		Vậy $M = \dfrac{28}{9}$.
	\end{enumerate}
}
\end{bt}

\begin{bt}%[9D8H26-3]
% nguon: Đề tham khảo Tuyển sinh 10 TP.HCM 2026-2027 – Cụm 2 (Đề 3), Bài 3, trang 13
Một hộp kín chứa $5$ loại thẻ có ghi tên lần lượt là $5$ từ THÀNH, PHỐ, HỒ, CHÍ, MINH với số lượng tương ứng là: THÀNH: $5$ thẻ; PHỐ: $4$ thẻ; HỒ: $7$ thẻ; CHÍ: $2$ thẻ; MINH: $5$ thẻ. Sau khi trộn đều các thẻ trong hộp, một học sinh rút ngẫu nhiên $1$ tấm thẻ.
\begin{enumerate}[a)]
	\item Trong hộp có bao nhiêu thẻ loại tên ngắn (tên thẻ chứa không quá $3$ chữ cái)?
	\item Tính xác suất của biến cố $A$: "Học sinh rút được loại thẻ không phải tên ngắn".
\end{enumerate}
\loigiai{
	\begin{enumerate}[a)]
		\item Các thẻ loại tên ngắn (chứa không quá $3$ chữ cái) gồm:
		\begin{itemize}
			\item Từ PHỐ ($3$ chữ cái): $4$ thẻ.
			\item Từ HỒ ($2$ chữ cái): $7$ thẻ.
			\item Từ CHÍ ($3$ chữ cái): $2$ thẻ.
		\end{itemize}
		Số lượng thẻ loại tên ngắn có trong hộp là:
		$$4 + 7 + 2 = 13\text{ (thẻ)}.$$
		\item Tổng số thẻ có trong hộp là:
		$$n(\Omega) = 5 + 4 + 7 + 2 + 5 = 23\text{ (thẻ)}.$$
		Biến cố $A$: "Học sinh rút được loại thẻ không phải tên ngắn", tức là thẻ rút được có tên chứa từ $4$ chữ cái trở lên.\\
		Các loại thẻ không phải tên ngắn gồm:
		\begin{itemize}
			\item Từ THÀNH ($5$ chữ cái): $5$ thẻ.
			\item Từ MINH ($4$ chữ cái): $5$ thẻ.
		\end{itemize}
		Số kết quả thuận lợi cho biến cố $A$ là:
		$$n(A) = 5 + 5 = 10\text{ (thẻ)}.$$
		Xác suất của biến cố $A$ là:
		$$P(A) = \frac{n(A)}{n(\Omega)} = \frac{10}{23}.$$
	\end{enumerate}
}
\end{bt}

\begin{bt}%[9D6V21-3]
% nguon: Đề tham khảo Tuyển sinh 10 TP.HCM 2026-2027 – Cụm 2 (Đề 3), Bài 4, trang 13
Một khu đất hình chữ nhật có chiều rộng $35\text{ m}$ được chia thành hai khu vườn nhỏ để trồng rau. Xung quanh hai khu vườn người ta làm lối đi: lối đi giữa hai vườn rau rộng $1{,}5\text{ m}$ và các lối đi còn lại quanh khu đất rộng $0{,}5\text{ m}$. Khu vườn thứ hai có chiều dài ít hơn khu vườn thứ nhất là $15\text{ m}$. Gọi $x\text{ (m)}$ là chiều dài của khu vườn thứ nhất ($x > 15$).
\begin{enumerate}[a)]
	\item Viết biểu thức biểu diễn theo $x$ tổng diện tích trồng rau của hai khu vườn.
	\item Tìm diện tích của hai khu vườn trồng rau, biết diện tích toàn bộ khu đất lớn hơn tổng diện tích trồng rau là $162{,}5\text{ m}^2$.
\end{enumerate}
\loigiai{
	\begin{enumerate}[a)]
		\item Chiều rộng của mỗi khu vườn rau sau khi bớt lối đi hai bên là:
		$$35 - 0{,}5 - 0{,}5 = 34\text{ (m)}.$$
		Khu vườn rau thứ nhất có chiều dài $x\text{ (m)}$ nên diện tích là $34x\text{ (m}^2\text{)}$.\\
		Khu vườn rau thứ hai có chiều dài là $x - 15\text{ (m)}$ nên diện tích là $34(x - 15)\text{ (m}^2\text{)}$.\\
		Tổng diện tích trồng rau của hai khu vườn là:
		$$S = 34x + 34(x - 15) = 34x + 34x - 510 = 68x - 510\text{ (m}^2\text{)}.$$
		\item Chiều dài tổng thể của toàn bộ khu đất (gồm cả các lối đi) là:
		$$0{,}5 + x + 1{,}5 + (x - 15) + 0{,}5 = 2x - 12{,}5\text{ (m)}.$$
		Diện tích của toàn bộ khu đất hình chữ nhật là:
		$$S_{\text{khu đất}} = 35(2x - 12{,}5) = 70x - 437{,}5\text{ (m}^2\text{)}.$$
		Theo giả thiết, diện tích khu đất lớn hơn diện tích trồng rau là $162{,}5\text{ m}^2$, ta có phương trình:
		$$(70x - 437{,}5) - (68x - 510) = 162{,}5 \iff 2x + 72{,}5 = 162{,}5 \iff 2x = 90 \iff x = 45\text{ (nhận).}$$
		Vậy tổng diện tích trồng rau của cả hai khu vườn là:
		$$S = 68 \cdot 45 - 510 = 3060 - 510 = 2550\text{ (m}^2\text{)}.$$
	\end{enumerate}
}
\end{bt}

\begin{bt}%[9H0V31-3]
% nguon: Đề tham khảo Tuyển sinh 10 TP.HCM 2026-2027 – Cụm 2 (Đề 3), Bài 5, trang 14
Một bồn chứa nước hình trụ có đường kính đáy là $1{,}4\text{ m}$ và chiều cao $3{,}25\text{ m}$. Người ta đổ nước vào bồn sao cho chiều cao của cột nước bằng một nửa chiều cao của bồn, sau đó thả vào một phao hình cầu bằng kim loại đặc không thấm nước có bán kính $30\text{ cm}$ và phao chìm hoàn toàn trong nước.
\begin{enumerate}[a)]
	\item Tính thể tích nước có trong bồn (làm tròn kết quả đến hàng phần mười).
	\item Sau đó người ta tiếp tục bơm thêm nước vào bồn bằng một máy bơm có công suất $0{,}0024\text{ m}^3/\text{giây}$. Hỏi sau bao nhiêu phút thì bồn đầy nước? (Làm tròn kết quả đến hàng đơn vị). Biết công thức thể tích hình trụ là $V = \pi r^2 h$, thể tích hình cầu là $V = \dfrac{4}{3}\pi R^3$.
\end{enumerate}
\loigiai{
	\begin{enumerate}[a)]
		\item Bán kính đáy của bồn hình trụ là:
		$$R = \frac{1{,}4}{2} = 0{,}7\text{ (m)}.$$
		Đổi bán kính phao hình cầu: $r = 30\text{ cm} = 0{,}3\text{ m}$.\\
		Thể tích phần hình trụ chứa nước và phao (chiều cao bằng một nửa bồn) là:
		$$V_1 = \pi R^2 \cdot \frac{h}{2} = \pi \cdot 0{,}7^2 \cdot \frac{3{,}25}{2} = 0{,}79625\pi\text{ (m}^3\text{)}.$$
		Thể tích của phao hình cầu là:
		$$V_{\text{phao}} = \frac{4}{3}\pi r^3 = \frac{4}{3}\pi \cdot 0{,}3^3 = 0{,}036\pi\text{ (m}^3\text{)}.$$
		Thể tích nước thực tế có trong bồn là:
		$$V_{\text{nước}} = V_1 - V_{\text{phao}} = 0{,}79625\pi - 0{,}036\pi = 0{,}76025\pi \approx 2{,}4\text{ (m}^3\text{)}.$$
		\item Vì mực chất lỏng (nước và phao) đang chiếm đúng một nửa thể tích của bồn, nên thể tích phần không gian còn trống cần bơm thêm nước để đầy bồn chính là:
		$$V_{\text{trống}} = \pi R^2 \cdot \frac{h}{2} = 0{,}79625\pi\text{ (m}^3\text{)}.$$
		Thời gian máy bơm cần chạy để làm đầy bồn là:
		$$t = \frac{V_{\text{trống}}}{0{,}0024} = \frac{0{,}79625\pi}{0{,}0024} \approx 1042{,}29\text{ (giây)}.$$
		Đổi sang phút:
		$$\frac{1042{,}29}{60} \approx 17{,}37\text{ (phút)} \approx 17\text{ (phút)}.$$
		Vậy sau khoảng $17$ phút thì bồn đầy nước.
	\end{enumerate}
}
\end{bt}

\begin{bt}%[9D1H3-4]
% nguon: Đề tham khảo Tuyển sinh 10 TP.HCM 2026-2027 – Cụm 2 (Đề 3), Bài 6, trang 14
Anh Bình là công nhân của một khu chế xuất. Trong tháng 5 vừa qua, hệ thống chấm công ghi nhận anh Bình đã làm việc tổng cộng $212$ giờ gồm các giờ làm trong định mức và giờ làm thêm ngoài định mức. Tiền công mỗi giờ làm trong định mức là $38\,000$ đồng, còn mỗi giờ làm thêm ngoài định mức được trả bằng $150\%$ tiền công một giờ trong định mức. Cuối tháng 5, anh Bình nhận được tổng số tiền lương là $8\,436\,000$ đồng. Hỏi trong tháng 5 anh Bình đã làm thêm bao nhiêu giờ ngoài định mức?
\loigiai{
	Gọi số giờ làm việc trong định mức là $x$ và số giờ làm thêm ngoài định mức là $y$ ($0 < x, y < 212$).\\
	Tổng số giờ làm việc trong tháng 5 là $212$ giờ, ta có phương trình:
	$$x + y = 212 \quad (1)$$
	Tiền công cho mỗi giờ làm thêm ngoài định mức là:
	$$38\,000 \cdot 150\% = 57\,000\text{ (đồng)}.$$
	Tổng tiền lương anh Bình nhận được là $8\,436\,000$ đồng, ta có phương trình:
	$$38\,000x + 57\,000y = 8\,436\,000 \quad (2)$$
	Từ $(1)$ và $(2)$, ta có hệ phương trình:
	$$\begin{cases} x + y = 212 \\ 38\,000x + 57\,000y = 8\,436\,000 \end{cases} \iff \begin{cases} 38x + 38y = 8056 \\ 38x + 57y = 8436 \end{cases} \iff \begin{cases} 19y = 380 \\ x = 212 - y \end{cases} \iff \begin{cases} y = 20 \\ x = 192 \end{cases}\text{ (thỏa mãn).}$$
	Vậy trong tháng 5 anh Bình đã làm thêm $20$ giờ ngoài định mức.
}
\end{bt}

\begin{bt}%[9H9C29-3]
% nguon: Đề tham khảo Tuyển sinh 10 TP.HCM 2026-2027 – Cụm 2 (Đề 3), Bài 7, trang 14
Cho tam giác $ABC$ có ba góc nhọn ($AB < AC$) nội tiếp đường tròn $(O; R)$. Hai đường cao $AD, BE$ của tam giác $ABC$ cắt nhau tại trực tâm $H$. Gọi $M, N$ lần lượt là giao điểm của đường tròn $(O)$ với các tia $BE, AD$ ($M \neq B, N \neq A$).
\begin{enumerate}[a)]
	\item Chứng minh tứ giác $ABDE$ nội tiếp đường tròn và xác định tâm $I$ của đường tròn này, từ đó suy ra $DE \parallel MN$.
	\item Kẻ đường kính $CK$ của đường tròn $(O)$. Chứng minh tứ giác $AKBH$ là hình bình hành và suy ra ba điểm $H, I, K$ thẳng hàng.
	\item Trong trường hợp $\widehat{BCA} = 60^\circ$, chứng minh $DE = \dfrac{1}{2}AB$ và tính diện tích hình viên phân giới hạn bởi cung nhỏ $DE$ và dây cung $DE$ của đường tròn $(I)$ theo $R$.
\end{enumerate}
\loigiai{
	\begin{center}
	\begin{tikzpicture}[scale=1]
		\coordinate (O) at (0,0);
		\def\R{3}
		\draw[thick] (O) circle (\R);
		\coordinate (A) at (75:\R);
		\coordinate (B) at (200:\R);
		\coordinate (C) at (-20:\R);
		
		\coordinate (D) at ($(B)!(A)!(C)$);
		\coordinate (E) at ($(A)!(B)!(C)$);
		\coordinate (H) at (intersection of A--D and B--E);
		
		\coordinate (K) at ($(O)!-1!(C)$);
		\coordinate (I) at ($(A)!0.5!(B)$);
		
		\coordinate (M) at (105:\R);
		\coordinate (N) at (275:\R);
		
		\draw[thick] (A) -- (B) -- (C) -- cycle;
		\draw[thick, blue] (A) -- (D);
		\draw[thick, blue] (B) -- (E);
		\draw[dashed] (D) -- (N);
		\draw[dashed] (E) -- (M);
		\draw[thick, red] (D) -- (E);
		\draw[dashed] (M) -- (N);
		\draw[thick] (C) -- (K);
		\draw[dashed] (A) -- (K) -- (B);
		\draw[dashed] (H) -- (K);
		
		\fill (A) circle (1.5pt) node[above] {$A$};
		\fill (B) circle (1.5pt) node[below left] {$B$};
		\fill (C) circle (1.5pt) node[below right] {$C$};
		\fill (D) circle (1.5pt) node[below] {$D$};
		\fill (E) circle (1.5pt) node[above right] {$E$};
		\fill (H) circle (1.5pt) node[below left] {$H$};
		\fill (O) circle (1.5pt) node[right] {$O$};
		\fill (K) circle (1.5pt) node[above left] {$K$};
		\fill (I) circle (1.5pt) node[above left] {$I$};
		\fill (M) circle (1.5pt) node[above left] {$M$};
		\fill (N) circle (1.5pt) node[below] {$N$};
	\end{tikzpicture}
	\end{center}
	\begin{enumerate}[a)]
		\item Xét tam giác $AEB$ vuông tại $E$ (do $BE \perp AC$) nên tam giác $AEB$ nội tiếp đường tròn đường kính $AB$.\\
		Xét tam giác $ADB$ vuông tại $D$ (do $AD \perp BC$) nên tam giác $ADB$ nội tiếp đường tròn đường kính $AB$.\\
		Do đó bốn điểm $A, B, D, E$ cùng thuộc đường tròn đường kính $AB$, suy ra tứ giác $ABDE$ nội tiếp đường tròn có tâm $I$ là trung điểm của $AB$.\\
		Trong tứ giác nội tiếp $ABDE$, ta có:
		$$\widehat{BDE} = \widehat{BAE} \quad (\text{hai góc nội tiếp cùng chắn cung } BE).$$
		Trong đường tròn $(O)$, ta có:
		$$\widehat{BMN} = \widehat{BAN} \quad (\text{hai góc nội tiếp cùng chắn cung } BN).$$
		Mà điểm $E$ nằm trên cạnh $AC$ và điểm $N$ nằm trên tia $AD$ nên $\widehat{BAE} = \widehat{BAC} = \widehat{BAN}$.\\
		Suy ra $\widehat{BDE} = \widehat{BMN}$.\\
		Hai góc này ở vị trí đồng vị nên $DE \parallel MN$.
		\item Ta có $\widehat{KAC} = 90^\circ$ (góc nội tiếp chắn nửa đường tròn đường kính $KC$), suy ra $KA \perp AC$.\\
		Mà $BE \perp AC$ (do $BE$ là đường cao) $\implies KA \parallel BH$.\\
		Tương tự, $\widehat{KBC} = 90^\circ$ (góc nội tiếp chắn nửa đường tròn đường kính $KC$), suy ra $KB \perp BC$.\\
		Mà $AD \perp BC$ (do $AD$ là đường cao) $\implies KB \parallel AH$.\\
		Tứ giác $AKBH$ có các cặp cạnh đối song song nên là hình bình hành.\\
		Vì $I$ là trung điểm của đường chéo $AB$ nên $I$ cũng là trung điểm của đường chéo $HK$.\\
		Vậy ba điểm $H, I, K$ thẳng hàng.
		\item Xét $\Delta CDE$ và $\Delta CAB$ có:
		$$\widehat{C}\text{ là góc chung và } \widehat{CDE} = \widehat{CAB}\text{ (cùng bù với } \widehat{BDE}\text{)}.$$
		Suy ra $\Delta CDE \backsim \Delta CAB$ (g-g), do đó:
		$$\frac{DE}{AB} = \frac{CD}{CA} = \cos \widehat{BCA} = \cos 60^\circ = \frac{1}{2} \implies DE = \frac{1}{2}AB.$$
		Kẻ đường kính $AQ$ của đường tròn $(O)$.\\
		Ta có $\widehat{ABQ} = 90^\circ$ và $\widehat{AQB} = \widehat{ACB} = 60^\circ$ (hai góc nội tiếp cùng chắn cung $AB$).\\
		Xét tam giác vuông $ABQ$ vuông tại $B$:
		$$AB = AQ \cdot \sin \widehat{AQB} = 2R \cdot \sin 60^\circ = R\sqrt{3}.$$
		Suy ra bán kính của đường tròn $(I)$ là:
		$$R_I = \frac{AB}{2} = \frac{R\sqrt{3}}{2}.$$
		Tam giác $IDE$ có $ID = IE = R_I$ và $DE = \dfrac{1}{2}AB = R_I$, do đó $\Delta IDE$ là tam giác đều cạnh $R_I = \dfrac{R\sqrt{3}}{2}$, suy ra $\widehat{DIE} = 60^\circ$.\\
		Diện tích hình quạt tròn $IDE$ là:
		$$S_{\text{quạt}} = \frac{\pi R_I^2 \cdot 60^\circ}{360^\circ} = \frac{\pi \cdot \left(\frac{R\sqrt{3}}{2}\right)^2}{6} = \frac{\pi R^2}{8}.$$
		Diện tích tam giác đều $IDE$ là:
		$$S_{\Delta IDE} = \frac{R_I^2 \sqrt{3}}{4} = \frac{\left(\frac{R\sqrt{3}}{2}\right)^2 \sqrt{3}}{4} = \frac{3R^2\sqrt{3}}{16}.$$
		Diện tích hình viên phân giới hạn bởi cung nhỏ $DE$ và dây cung $DE$ là:
		$$S = S_{\text{quạt}} - S_{\Delta IDE} = \frac{\pi R^2}{8} - \frac{3R^2\sqrt{3}}{16} = \frac{(2\pi - 3\sqrt{3})R^2}{16}.$$
	\end{enumerate}
}
\end{bt}
